\documentclass[../../../include/open-logic-section]{subfiles}

\begin{document}

\olfileid{sth}{z}{story}
\olsection{الرواية بمزيد من التفصيل}
	
اقتبسنا في~\olref[story][approach]{sec} وصف شونفيلد لعملية تكوين
المجموعات. ونريد الآن تدوين بضعة مبادئ أخرى لكي نجعل هذه الرواية أدق
قليلًا. وهذه هي:
\begin{enumerate}
\item[] \stageshier. تتكون كل مجموعة في مرحلة ما.
\item[] \stagesord. المراحل مرتبة: فبعضها يأتي \emph{قبل}
بعض.\footnote{سنفترض في الواقع، ضمنًا، أن المراحل \emph{مرتبة
  ترتيبًا جيدًا}. ويُشرح معنى هذا في~\olref[ordinals][]{chap}.
  وهذا افتراض جوهري. بل يمكن، باستعمال تقنية شديدة الذكاء ترجع إلى
  \citet{Scott1974}، \emph{تفادي} هذا الافتراض ثم \emph{اشتقاقه}.
  (وسيشرح هذا أيضًا لماذا ينبغي أن نعتقد بوجود مرحلة ابتدائية.) لا
  يمكننا الخوض في ذلك هنا؛ للمزيد، انظر~\citet{ButtonLT1}.}
\item[] \stagesacc. لكل مرحلة~$S$، ولكل مجموعات تكونت \emph{قبل}
المرحلة~$S$: تتكون في المرحلة~$S$ مجموعة عناصرها تلك المجموعات نفسها
دون غيرها. ولا يتكون شيء آخر في المرحلة~$S$.
\end{enumerate}
هذه مبادئ غير صورية، لكننا سنستطيع استعمالها لتسويغ عدد من مسلّمات
نظرية زرملو للمجموعات.

(ينبغي أن نقدم كلمة تحذير. فرغم أننا سنعرض مسلّمات معيارية تمامًا
بأسماء معيارية تمامًا، فإن المبادئ المكتوبة بخط مائل التي قدمناها للتو
ليس لها أسماء معينة في الأدبيات. لقد أعطيناها ببساطة ألقابًا نأمل أن
تكون نافعة.)

\end{document}
